% !TEX program = xelatex
\documentclass[11pt,a4paper]{article}
\usepackage{fontspec}
\setmainfont{Calibri}
\setmonofont[Scale=0.88]{Consolas}
\newfontfamily\symbolfont{Segoe UI Symbol}
\usepackage{polyglossia}\setdefaultlanguage{polish}
\usepackage[table]{xcolor}
\usepackage{booktabs,longtable,array,geometry,enumitem,graphicx,fancyhdr,caption}
\usepackage[most]{tcolorbox}
\PassOptionsToPackage{hyphens}{url}
\usepackage[hidelinks,unicode]{hyperref}
\emergencystretch=2.5em
\geometry{a4paper,margin=2.1cm,top=2.3cm,bottom=2.3cm}
\setlength{\parindent}{0pt}\setlength{\parskip}{5pt}
\definecolor{apteal}{HTML}{0E7C86}
\definecolor{apteallight}{HTML}{E6F4F5}
\definecolor{apgrey}{HTML}{F3F4F6}
\definecolor{apamber}{HTML}{B45309}
\definecolor{apamberlight}{HTML}{FEF3C7}
\definecolor{apred}{HTML}{B91C1C}
\definecolor{apredlight}{HTML}{FEE2E2}
\definecolor{apgreen}{HTML}{15803D}
\definecolor{apgreenlight}{HTML}{DCFCE7}
\renewcommand{\arraystretch}{1.25}
\captionsetup{font=small,labelfont=bf}
\pagestyle{fancy}\fancyhf{}
\fancyfoot[L]{\small CZYTAJ.md}
\fancyfoot[R]{\small\thepage}
\renewcommand{\headrulewidth}{0pt}
\newtcolorbox{uwaga}[1][Uwaga]{enhanced,breakable,colback=apamberlight,colframe=apamber,title={#1},fonttitle=\bfseries,boxrule=0.6pt,arc=2pt,left=6pt,right=6pt,top=4pt,bottom=4pt}
\newtcolorbox{info}[1][Jak to działa]{enhanced,breakable,colback=apteallight,colframe=apteal,title={#1},fonttitle=\bfseries,boxrule=0.6pt,arc=2pt,left=6pt,right=6pt,top=4pt,bottom=4pt}
\newtcolorbox{wazne}[1][Ważne]{enhanced,breakable,colback=apredlight,colframe=apred,title={#1},fonttitle=\bfseries,boxrule=0.6pt,arc=2pt,left=6pt,right=6pt,top=4pt,bottom=4pt}
\newtcolorbox{przyklad}[1][Przykład liczbowy]{enhanced,breakable,colback=apgreenlight,colframe=apgreen,title={#1},fonttitle=\bfseries,boxrule=0.6pt,arc=2pt,left=6pt,right=6pt,top=4pt,bottom=4pt}
\newtcolorbox{kodbox}{enhanced,breakable,colback=apgrey,colframe=apgrey,boxrule=0pt,arc=2pt,left=5pt,right=5pt,top=3pt,bottom=3pt,fontupper=\ttfamily\footnotesize}
\setlist{nosep,leftmargin=*}
\setcounter{secnumdepth}{2}
\begin{document}

{\LARGE\bfseries Q11 w Home Assistant: łatka integracji Shelly\par}
\vspace{4pt}\textcolor{apteal}{\rule{\textwidth}{1.2pt}}

Autor: Dawid Cekała, 29.09.2026


\section*{Po co}

Oficjalna integracja Shelly w Home Assistant tworzy encje pól Shelly X tylko wtedy, gdy klucz pola jest na liście w jej kodzie. Z 15 pól Q11 na tej liście są tylko fazy. Limit prądu i włącznik ładowania są tam zapisane wyłącznie dla ładowarki TopAC (kod produktu EVE01). Bez łatki Home Assistant pokazuje 19 encji: 12 fazowych i 7 systemowych. Z łatką pokazuje 33 encje.


\section*{Co zmienia łatka}

\begin{longtable}{>{\raggedright\arraybackslash}p{2.35cm}>{\raggedright\arraybackslash}p{13.37cm}}
\toprule
\rowcolor{apteallight}\textbf{Plik} & \textbf{Zmiana} \\
\midrule\endhead
\bottomrule\endfoot
const.py & lista kodów produktu Ampere Point, dziś \texttt{apq11dev} \\
entity.py & pole wyboru bez napisów opcji nie wywraca encji (dotyczy trybu pracy i sygnału pojazdu) \\
number.py & limit prądu także dla Q11; nowe: limit energii, okno ładowania od i do \\
switch.py & włącznik ładowania \\
select.py & tryb ładowania (od razu, do limitu energii, w oknie godzin) \\
sensor.py & stan ładowarki, sygnał pojazdu, energia i czas sesji, energia ostatniej sesji, temperatura sterownika, usterki, wersja sterownika \\
translations & nazwy po polsku i angielsku \\
\end{longtable}

Pełna różnica względem Home Assistant 2026.6.4: \texttt{latka\_\allowbreak{}shelly\_\allowbreak{}q11\_\allowbreak{}HA2026.\allowbreak{}6.\allowbreak{}4.\allowbreak{}diff}.


\section*{Jak działa}

Home Assistant najpierw szuka integracji w \texttt{C:\allowbreak{}\textbackslash{}HomeAssistant\textbackslash{}config\textbackslash{}custom\_\allowbreak{}components}. Kopia integracji Shelly w tym folderze zastępuje wbudowaną. Dotyczy to wszystkich urządzeń Shelly w tym Home Assistant, ale inne urządzenia działają jak dotąd, bo łatka dopisuje wyłącznie encje dla kodu \texttt{apq11dev}.


\section*{Po aktualizacji Home Assistant}

Kopia zostaje w wersji 2026.6.4 i po aktualizacji może przestać pasować do nowego Home Assistant. Wtedy integracja Shelly się nie uruchomi. Po każdej aktualizacji:


\begin{kodbox}
\mbox{}docker\ exec\ homeassistant\ rm\ -rf\ /config/custom\_components/shelly\\
\mbox{}docker\ exec\ homeassistant\ cp\ -r\ /usr/src/homeassistant/homeassistant/components/shelly\ /config/custom\_co\\
\hspace*{1.2em}mponents/shelly\\
\mbox{}python\ latka\_shelly\_q11.py\ C:\textbackslash{}HomeAssistant\textbackslash{}config\textbackslash{}custom\_components\textbackslash{}shelly\\
\mbox{}docker\ restart\ homeassistant
\end{kodbox}

Jeśli łatka zgłosi błąd (zmienił się kod integracji), wystarczy usunąć folder i zrestartować Home Assistant. Wróci wtedy stan sprzed łatki z 19 encjami.


\section*{Docelowo}

Ta sama zmiana może trafić do oficjalnej integracji. Tak dopisano ładowarkę TopAC. Wymaga to ostatecznego kodu produktu zamiast \texttt{apq11dev}.


\end{document}